\documentclass[12pt,a4paper]{report}

% Paquetes esenciales
\usepackage[utf8]{inputenc}
\usepackage[spanish]{babel} % Idioma español para "Índice", "Capítulo", etc.
\usepackage{geometry}
\geometry{a4paper, margin=2.5cm}
\usepackage{fancyhdr} % Para personalizar encabezados y pies de página
\usepackage{hyperref} % Para que el índice tenga enlaces clickeables
\usepackage{xcolor} % Para darle color a los enlaces
\usepackage{titlesec} % Para dar formato a los títulos de capítulos

% Configuración de colores para los enlaces
\hypersetup{
    colorlinks=true,
    linkcolor=blue,
    filecolor=magenta,      
    urlcolor=cyan,
    pdftitle={FAQ Cachimbos FIM 2026-1},
}

% Configuración de Encabezados
\pagestyle{fancy}
\fancyhf{}
\fancyhead[L]{\textit{\leftmark}} % Muestra el nombre del capítulo a la izquierda
\fancyhead[R]{\thepage} % Muestra el número de página a la derecha
\renewcommand{\headrulewidth}{0.4pt} % Línea divisoria en el encabezado

% Formato de los capítulos para que no ocupen media página en blanco
\titleformat{\chapter}[display]
  {\normalfont\bfseries\Huge}{\chaptertitlename\ \thechapter}{20pt}{\Huge}
\titlespacing*{\chapter}{0pt}{-20pt}{20pt}

\begin{document}

% ---------------- CARÁTULA ----------------
\begin{titlepage}
    \centering
    \vspace*{2cm}
    
    {\Huge \textbf{Manual de Supervivencia}}\\[0.5cm]
    {\Huge \textbf{FIM 2026-1}}\\[1.5cm]
    
    {\Large Guía Definitiva de Trámites, Pagos y Vida Universitaria}\\[2cm]
    
    {\Large \textbf{Autor:} Gustavo Morales}\\[1cm] % Personalizado con tu nombre
    
    {\large Facultad de Ingeniería Mecánica}\\[0.5cm]
    {\large Universidad Nacional de Ingeniería}\\[3cm]
    
    \vfill
    
    \hrule
    \vspace{0.5cm}
    \textbf{¿Qué encontrarás en este documento?}\\[0.5cm]
    \begin{flushleft}
    Este manual ha sido recopilado para resolver las dudas más frecuentes de los nuevos ingresantes de la FIM. Aquí detallamos paso a paso cómo lidiar con la saturación de la plataforma INTRALU, los pagos de OCEF, la toma de fotografías de DIRCE, el examen médico, y te damos los consejos esenciales para sobrevivir a la carga académica y entender la cultura de la facultad sin perder la paciencia.
    \end{flushleft}
    \vspace{0.5cm}
    \hrule
\end{titlepage}

% ---------------- ÍNDICE ----------------
\tableofcontents
\newpage

% ---------------- CAPÍTULO 1 ----------------
\chapter{Trámites y Pagos (INTRALU)}

\section{Acceso a la Plataforma}
El sistema suele saturarse debido a la gran cantidad de ingresantes. Si tienes problemas para acceder a tu INTRALU, asegúrate de iniciar sesión mediante la opción de Google utilizando tu correo institucional (@uni.pe).

\section{Periodos Académicos y Generación de Órdenes}
Una de las mayores confusiones ocurre con los periodos de pago. Aunque tu ingreso es para el ciclo \textbf{2026-1}, el periodo administrativo actual en el sistema sigue siendo \textbf{2025-3}. 
Si intentas generar el pago en 2026-1 y no te lo permite, o si te sale "orden no registrada", debes generar tus pagos seleccionándolos manualmente de la lista de todos los trámites.

\section{¿Qué pagar primero?}
No es necesario pagar todo de golpe. Prioriza únicamente:
\begin{itemize}
    \item \textbf{Examen Médico:} Tiene un costo aproximado de S/ 82.00.
    \item \textbf{Autoseguro}
\end{itemize}
Los pagos para el carné universitario y el carné de biblioteca pueden esperar a que inicie oficialmente el ciclo regular. Una vez generada la orden en INTRALU, puedes pagar tranquilamente a través de la aplicación del BCP o Yape.

% ---------------- CAPÍTULO 2 ----------------
\chapter{Fotografías DIRCE y Constancias}

\section{Demora en la actualización de fotos}
Luego de entregar tus documentos presencialmente, te tomarán una fotografía. Es completamente normal que esta imagen no aparezca de inmediato en tu INTRALU; al haber cientos de ingresantes, DIRCE procesa esta información manualmente y puede tardar varios días, o incluso hasta el inicio del ciclo.

\section{Constancia de Ingreso}
La foto es un requisito indispensable para que el sistema te permita descargar la constancia de ingreso final. Debes tener paciencia y revisar el portal periódicamente.

% ---------------- CAPÍTULO 3 ----------------
\chapter{Examen Médico}

El examen médico es un paso obligatorio. Para poder ser atendido ese día, es indispensable que presentes tu voucher o comprobante de pago (puedes llevarlo impreso o mostrarlo desde tu celular). 
Si no has pagado, no podrás pasar el examen. Ve bien aseado, de preferencia temprano para evitar largas colas, y lleva bloqueador solar y agua. Si usas brackets, no te preocupes, el personal simplemente lo anotará en tu registro.

% ---------------- CAPÍTULO 4 ----------------
\chapter{Ciclo Introductorio y Ciclo Regular}

\section{Clases Introductorias}
Están programadas para iniciar el lunes 2 de marzo. Si bien no son obligatorias y sus evaluaciones no afectan tu promedio, son altamente recomendables, especialmente para cursos filtro como Física I. Te enseñarán los mismos ingenieros que dictan en el ciclo regular.

\section{Inicio del Ciclo Regular}
El ciclo 2026-1 iniciará oficialmente el 23 de marzo y terminará aproximadamente la última semana de julio. No te preocupes por armar tu horario ahora; la primera matrícula para los cachimbos es completamente automática.

% ---------------- CAPÍTULO 5 ----------------
\chapter{Mitos y Vida en la FIM}

\section{Rendimiento Académico: Bikas y Trikas}
En la UNI, desaprobar tiene consecuencias que debes evitar a toda costa:
\begin{itemize}
    \item \textbf{Bika (Jalar 2 veces un curso):} Te consideran alumno en riesgo académico. Recibirás tutoría obligatoria, atención psicológica y, lo peor, se limitará la cantidad de créditos/cursos que puedes llevar el siguiente ciclo, atrasándote considerablemente.
    \item \textbf{Trika (Jalar 3 veces un curso):} Separación automática de la universidad por un año. Si al regresar vuelves a desaprobar ese mismo curso, la expulsión es permanente.
\end{itemize}

\section{Centros de Estudiantes y Agrupaciones}
La facultad fomenta la investigación y el desarrollo extracurricular. Te recomendamos acercarte a los centros de estudiantes según tu especialidad:
\begin{itemize}
    \item \textbf{CIIM:} Ingeniería Mecánica (M3)
    \item \textbf{CEDIME:} Ingeniería Mecánica Eléctrica (M4)
    \item \textbf{APEIN:} Ingeniería Naval (M5)
    \item \textbf{CEDIM:} Ingeniería Mecatrónica (M6)
    \item \textbf{ACIA:} Asociación Cultural Isaac Asimov. Ideal si buscas clubes de música, cine, fotografía, lectura, arte o idiomas.
\end{itemize}

\section{El Mito del Carné de Comedor}
No existe tal cosa como un "Carné de comedor". Si un alumno mayor te ofrece tramitarlo o te pide dinero para agilizarlo, te está engañando (kchimbeando).

\end{document}
